\begin{abstract}
Large vision--language models (LVLMs) frequently hallucinate, describing objects that are absent from the image and undermining their reliability in practical applications. Recent contrastive decoding methods mitigate this problem by comparing the model's standard output against a distorted counterpart, but they typically require multiple forward passes or complex training procedures. The ONLY method~\cite{wan2025only} achieves strong results with a single forward pass by suppressing attention heads with low text-to-visual entropy ratio at a single layer. However, this one-shot mask is frozen after layer~0, ignoring potentially valuable suppression evidence from deeper layers. We introduce Cumulative Mask Decoding (CMD), a training-free extension that computes a fresh head-suppression mask at every Transformer layer and blends it into a running cumulative mask via exponential moving average with a dynamic smoothing rate. This allows the suppression signal to be refined iteratively as the model processes deeper representations, while the adaptive EMA rate accelerates convergence when many heads exhibit imbalanced attention and dampens noise when few do. On the POPE benchmark with LLaVA-1.5, CMD achieves 88.00\% accuracy (+8.6 points over ONLY) on the challenging Adversarial split and matches ONLY on Popular, though at the cost of reduced precision on Random negatives. An analysis of the contrastive dynamics reveals that CMD operates almost entirely in the collaborative regime ($>98.6\%$ of decoding steps), indicating that the improvement stems from a more stable and reliable contrastive path rather than from more frequent contrastive corrections. CMD adds less than 0.5\% computational overhead, requiring no additional training or architectural modifications.
\end{abstract}